\section{Intrinsically Scoped, Extrinsically Typed Syntaxes}

\newcommand{\declentry}[2]{
    \AgdDta{#1}\AgdaSpace{}
  & \AgdaSymbol{:}\AgdaSpace{}
  & #2
}

\begin{frame}
  \frametitle{Mechanizing Type Theory: A Methodology}

  \pause I am working on an \textit{intrinsically scoped} and
  \textit{extrinsically typed} presentation of type
  theory\myfootcite{marquet}. The formal construction happens in two stages.
  \begin{tightenumerate}
    \item<+(1)-> Well-scoped, untyped syntax.
    \item<+(1)-> Typing derivations.
  \end{tightenumerate}

  \pause
  \begin{center}
    \begin{tabular}{l@{\hskip 3em}l@{\hskip 3em}l}
      \textit{Data Type} & \textit{Typing Derivation} & \textit{Meaning}
      \\
      \pause
      \\ \IsCtx{\variable{Γ}}
      &  \CtxTyping{\variable{Γ}}
      &  \variable{Γ} is a well-formed context
      \pause
      \\ \IsVar{\variable{v}}{\variable{Γ}}
      &  {\tiny \textit{(none)}}
      &  \variable{v} is a variable in \variable{Γ}
      \pause
      \\ \IsTyp{\variable{A}}{\variable{Γ}}
      &  \TypTyping{\variable{Γ}}{\variable{A}}
      &  \variable{A} is well-formed in context \variable{Γ}
      \pause
      \\ \IsTrm{\variable{t}}{\variable{Γ}}
      &  \TrmTyping{\variable{Γ}}{\variable{t}}{\variable{A}}
      &  \variable{t} has type \variable{A} in context \variable{Γ}
    \end{tabular}
  \end{center}

  \pause Specificity of my approach: notice how variables, types and terms
  are required to live in a context.
\end{frame}

\begin{frame}
  \frametitle{Russell/Tarski Universes}

  \pause In my work, I want types and terms to be distinct. There is a decoding
  operation that injects terms into types, together with the appropriate typing
  rules.

  \pause Gilbert et.~al.\myfootcite{gilbert} don't have this decoding
  operation as their types and terms are merged into expressions.

  \pause Merging types with terms allows for \textbf{Russell} universes,
  and keeping them separate forces a \textbf{Tarski} presentation.
\end{frame}

\begin{frame}
  \frametitle{Problem.}

  \pause Controlling relevance in Russell style:
  \begin{tightcenter}
    \begin{prooftree}
      \hypo{...}
      \hypo{Γ\,⊢\,A\,:\,\textsc{Type}}
      \hypo{...}
      \infer3{\text{$A$ is proof-relevant}}
    \end{prooftree}
    \hspace{3em}
    \begin{prooftree}
      \hypo{...}
      \hypo{Γ\,⊢\,A\,:\,\textsc{sProp}}
      \hypo{...}
      \infer3{\text{$A$ is proof-irrelevant}}
    \end{prooftree}
  \end{tightcenter}

  \pause \textbf{Problem.} Because I don't merge types and terms in
  expressions, I can't even write $A\,:\,\textsc{Prop}$ as above.

  \pause All I really want is an equivalent of the following rule
  \begin{tightcenter}
    \begin{prooftree}
      \hypo{Γ\,⊢\,A\,:\,\textsc{sProp}}
      \hypo{Γ\,⊢\,x\,:\,A}
      \hypo{Γ\,⊢\,y\,:\,A}
      \infer3{Γ\,⊢\,x\,≡\,y\,:\,A}
    \end{prooftree}
  \end{tightcenter}
\end{frame}

\begin{frame}
  \frametitle{Solution.}

  \pause All I really want is an equivalent of the following rule
  \begin{tightcenter}
    \begin{prooftree}
      \hypo{Γ\,⊢\,A\,:\,\textsc{sProp}}
      \hypo{Γ\,⊢\,x\,:\,A}
      \hypo{Γ\,⊢\,y\,:\,A}
      \infer3{Γ\,⊢\,x\,≡\,y\,:\,A}
    \end{prooftree}
  \end{tightcenter}

  \pause In this rule, all I really need to know if whether types are
  relevant or not.
  \begin{tightcenter}
    \textbf{Bake the relevance information into the syntax.}
  \end{tightcenter}
\end{frame}

\begin{frame}
  \frametitle{Baking relevance information into the syntax?}

  \pause For dependent type theory (without proof irrelevance), in my
  mechanization, I define contexts, variable, types and terms.

  \pause To ``bake relevance information into the syntax'', the idea is to
  split types and terms along the relevant/irrelevant distinction. We get
  \begin{tightitemize}
    \item<+(1)-> Contexs
    \item<+(1)-> Variables (relevant)
    \item<+(1)-> Variables (irrelevant)
    \item<+(1)-> Types (relevant)
    \item<+(1)-> Types (irrelevant)
    \item<+(1)-> Terms (relevant)
    \item<+(1)-> Terms (irrelevant)
  \end{tightitemize}

  \pause Sounds naïve... \\
  \textbf{Next section:} how this makes sense from the point of
  view of low-dimensional category theory!
\end{frame}

